% Folder 93, batch 11 — archivists' pages 201–220.
% Pass: Opus 5.5 (claude-opus-5-5), 2026-10-10 — first pass, unchecked against the pages by a human.
%
% Pages 201-203 close a run on the Cartier operation and Frobenius eigenvalues;
% 204, 213 and 218 are covers in his hand (their words become section headings);
% 205-212 are a numbered run on relative Picard functors with connections
% (P^c, P^ci, P^cip, P^cin); 214-216 on c^1(L) and connections; 215 and 217
% are cancelled sheets; 219-220 begin the canonical Gauss-Manin connection.
\documentclass[11pt,a4paper]{article}
\input{../preamble/grothendieck}

\folder{93}
\batch{11}
\pages{201}{220}
\foldertitle{Connexions de Gauss-Manin et équations de Picard-Fuchs. Opération de Cartier : copies de tapuscrit annoté (s.d.), tiré à part (1981), notes manuscrites (s.d.).}
\dating{[1972]-1981}
\watermark{Édition de démonstration}

\begin{document}

\page{201}
\note{l'argument vient d'avant le lot ; les trois lignes de la page sont encadrées et biffées de traits obliques}
\struck{Supposons que \uncertain{l'on ait} $S = \mathrm{Spec}\,\mathbb{F}_{p^{\nu}}$. Alors on \uncertain{peut} définir $F_{X}^{(\nu)}$, qui est un $S$-endomorphisme, et qui agit donc sur $\underline{\omega}$.}

\page{202}
Et on trouve donc que l'endomorphisme $F_{*}^{(\nu)}$ \struck{\ill{}} de $\underline{\omega}$ \add{$F^{(\nu)}_{X/S*}$} n'est autre que l'\emph{homom.\ transposé}, \uncertain{associé} à l'endomorphisme $F^{(\nu)}_{X/S} = (F_{X})^{\nu}$ de $X/S$, [\struck{qui} \uncertain{qui est} transposé de l'endomorphisme $F^{(\nu)\,*}_{X/S}$ sur $\mathbb{R}^{n-1}f_{*}(\underline{\Omega}^{n-1}_{X/S})$, \uncertain{grâce à} $X$ est \uncertain{propre} et lisse sur $S$\,…].

\struck{Plaçons-nous dans le \ill{}} Notons que nous avons la suite exacte
\[
0 \to \underline{\omega} \to \mathcal{H}^{1}_{\mathrm{DR}}(X/S) \to R^{1}f_{*}(\underline{\mathcal{O}}_{X/S})
\]
\note{sous $\underline{\omega}$ : « $= \mathcal{H}^{1,0}_{\mathrm{DR}}(X/S)$ » ; devant $R^{1}f_{*}$, un symbole biffé \ill{}}
[avec \uncertain{un} $0$ à droite dans les cas favorables), et que l'homom.\ de Cartier associé à $F^{(\nu)}_{X/S}$ \struck{est} compatible aux \uncertain{homom.\ de} $\equiv$ nous \uncertain{est} \ill{}
\[
\mathcal{H}^{1}_{\mathrm{DR}}(X/S) \to \mathcal{H}^{1}_{\mathrm{DR}}(X^{(p^{\nu}/S)}) \simeq \mathcal{H}^{1}_{\mathrm{DR}}(X/S)^{(p^{\nu})},
\]
\note{au-dessus du dernier $\simeq$, de sa main : « \uncertain{ss réserve} »}
(qui est un endomorphisme si $S = \mathrm{Spec}\,\mathbb{F}_{p^{\nu}}$). \struck{Donc}

D'ailleurs, dans ce cas \add{et si $X/S$ est propre}, \struck{\ill{}}, est \uncertain{bien} \uncertain{de} Cartier \uncertain{induit} zéro sur $\mathcal{H}^{1}_{\mathrm{DR}}(X/S)/\underline{\omega} \simeq R^{1}f_{*}(\underline{\mathcal{O}}_{X/S})$, \struck{il \uncertain{en} \uncertain{résulte} $\equiv$ \uncertain{une} \uncertain{sur} \uncertain{les}} (\uncertain{conjecture} \uncertain{de} \uncertain{transposition}), donc si $S = \mathrm{Spec}\,\mathbb{F}_{p^{\nu}}$, \uncertain{les} \struck{\uncertain{valeurs} qui sont $\neq 0$} \add{\uncertain{celles}} valeurs propres des \struck{\uncertain{valeurs}} (\ill{}) \add{\uncertain{celles} qui sont $\neq 0$} dans $\mathcal{H}^{1}_{\mathrm{DR}}(X/S)$ \uncertain{sont} \uncertain{aussi} \add{\uncertain{mod $p$}} les valeurs propres dans $\underline{\omega}$ et \uncertain{aussi} \uncertain{les} \uncertain{valeurs propres} de Frobenius $\ell$-adique ($\ell \neq p$) \struck{\uncertain{ou} cristallin} \uncertain{ou} cristalline qui sont des unités $p$-adiques (Katz).

On veut donc \uncertain{des} \uncertain{conditions} sur les $\omega$ telles que $\tilde{C}\omega = \tilde{\omega}$
\note{la phrase se poursuit en haut de la page 203}

\page{203}
\note{suite de la page 202 ; le développement se clôt ici (« Fin de parenthèse »)}
\ill{} qui sont loc.\ des diviseurs logarithmiques, permettant de construire des \struck{\uncertain{valeurs}} \add{\uncertain{valeurs}} propres de $F^{\nu}_{X/S}$ \struck{qui sont} (\struck{\ill{}} $\ell$-adiques, ou cristallins) qui sont $\equiv 1 \ (p)$.

Plus précisément, il existe $\omega$ avec $\tilde{C}\omega = \tilde{\omega}$, $\omega \neq 0$, ssi il existe une valeur propre de Frobenius congrue à 1 mod $p$. La dim de l'espace \struck{des} engendré par (sur $k$ !) les $\omega$ tels que $\tilde{C}\omega = \tilde{\omega}$ est égale au nb des valeurs propres de Frobenius congrues à 1 mod $p$. C'est aussi égal à la dim sur $\mathbb{F}_p$ du groupe \struck{\uncertain{des points de}} \struck{\uncertain{fonction}} ${}_{p}\mathbf{A}(k)$, $\mathbf{A}$ la jacobienne.

Fin de parenthèse\,..

\section{Variantes concernant des foncteurs de Picard relatifs}
\note{titre pris sur la page 204, une chemise portant de sa main, en haut à droite, « Variantes concernant des foncteurs de Picard relatifs », puis, souligné, « \uncertain{Att.}\ : »}

\page{205}
Soit $f : X \to S$ un morphisme de schémas, $f^{\mathrm{fl}}$ le morphisme correspondant des sites fppf.

① $P^{c} = P^{c}_{X/S} : (\mathrm{Sch})^{\circ}_{/S} \to (\mathrm{Ab})$, \struck{défini par} \add{faisceau associé au préfaisceau} $S' \mapsto$ classes de modules inversibles sur $X_{S'}$, munis d'une connexion rel.\ $/S'$.
\struck{$P^{c}(S') =$}
\[
(1.1) \qquad 0 \to f^{\mathrm{fl}}_{*}(\underline{\Omega}^{1}_{X/S})/\mathrm{Im}\, f^{\mathrm{fl}}_{*}(\underline{\mathcal{O}}^{*}_{X/S}) \to P^{c}_{X/S} \to \underline{\mathrm{Pic}}_{X/S} \to R^{1}f^{\mathrm{fl}}_{*}(\underline{\Omega}^{1}_{X/S})
\]
\marginal{$f^{\mathrm{fl}}_{*}(\underline{\mathcal{O}}^{*}_{X/S}) \to f^{\mathrm{fl}}_{*}(\underline{\Omega}^{1}_{X/S})$ donnée par $\varphi \mapsto d\varphi/\varphi$. C'est \ill{}, \uncertain{une} si \ill{} \uncertain{propre} et \ill{} \uncertain{à fibres géom. connexes}, \ill{} \uncertain{p.\ ex.\ propre} \ill{}}

\emph{NB} Si $f$ \struck{\uncertain{propre et plat}} \add{\uncertain{cohérent}}, et $\underline{\Omega}^{1}_{X/S}$ cohom.\ plat$/S$ en dim 0, ou $R^{1}f_{*}(\underline{\Omega}^{1}_{X/S})$ commute au chgt de base, alors pour tout schéma abélien $\mathbf{A}$ sur $S$, tout \uncertain{homom.}\ $B \to R^{1}f^{\mathrm{fl}}_{*}(\underline{\Omega}^{1}_{X/S})$ est nul. Donc si \ill{} \ill{} $u : B \to \underline{\mathrm{Pic}}_{X/S}$, on a $u(B) \subset \mathrm{Im}(P^{c} \to \underline{\mathrm{Pic}}_{X/S})$, d'où suite exacte
\[
(1.2) \qquad 0 \to f^{\mathrm{fl}}_{*}(\underline{\Omega}^{1}_{X/S}) \to u^{*}(P^{c}) \to B \to 1
\]
Ceci s'applique en particulier si $f$ propre sur $S$, et \ill{} \ill{}, et si $B = \underline{\mathrm{Pic}}^{0}_{X/S}$ est un schéma abélien. Alors, $u$ étant l'inclusion, on aura $u^{*}(P^{c}) = (P^{c})^{0}$, et on trouve suite exacte
\[
(1.3) \qquad 0 \to f^{\mathrm{fl}}_{*}(\underline{\Omega}^{1}_{X/S}) \to (P^{c}_{X/S})^{0} \to B \to 0, \qquad B = \underline{\mathrm{Pic}}^{0}_{X/S}
\]
\emph{Re NB} \ill{}, et s'il \ill{} \uncertain{lisse}, \uncertain{cas}, \ill{} fonctoriel en $X/S$, et commute au chgt de base.

② $P^{ci} = P^{ci}_{X/S}$, sous-groupe de $P^{c}_{X/S}$, faisceau associé au préfaisceau $C' \mapsto$ classes des modules inversibles sur $X_{S'}$, munis d'une connexion rel.\ à $S'$ qui est \emph{à courbure nulle}.

\page{206}
On a donc
\[
(2.1) \qquad 0 \to P^{ci}_{X/S} \to P^{c}_{X/S} \to f^{\mathrm{fl}}_{*}(Z^{2}_{X/S}) \subset f^{\mathrm{fl}}_{*}(\Omega^{2}_{X/S})
\]
\marginal{« courbure », avec un trait vers le terme $f^{\mathrm{fl}}_{*}(Z^{2}_{X/S})$}
ou, si on préfère, dans le cas où $X$ est lisse sur $S$,
\[
\begin{aligned}
(2.2) \qquad 0 &\to f^{\mathrm{fl}}_{*}(\underline{Z}^{1}_{X/S})/\mathrm{Im}\, f^{\mathrm{fl}}_{*}(\underline{\mathcal{O}}^{*}_{X/S}) \to P^{ci}_{X/S} \to \underline{\mathrm{Pic}}_{X/S} \\
&\to \mathbb{R}^{2}f^{\mathrm{fl}}_{*}(\mathcal{O} \to \underline{\Omega}^{1}_{X/S} \to \underline{\Omega}^{2}_{X/S} \to \cdots)
\end{aligned}
\]
[Dans le cas non lisse, on remplace le dernier terme par \struck{\ill{}} le faisceau associé : $S' \mapsto \mathbb{R}^{2}f_{S',\mathrm{cris}}(\mathcal{J}_{X_{S'}/S'})$ ← idéal d'augmentation de $(X_{S'}/S')_{\mathrm{cris}}$ ]
\note{au-dessus du passage biffé, « cris » de sa main}

\marginal{remis dans 1°)}
\emph{Cas particulier}. Plaçons-nous dans le cas où $\underline{\mathrm{Pic}}^{0}_{X/S}$ est un schéma abélien $B$, dont le dual est noté $A^{0} = \underline{\mathrm{Alb}}^{0}_{X/S}$. On a alors $X \to \underline{\mathrm{Alb}}^{1}_{X/S} = A^{1}$ (torseur sous $A^{0}$), et supposons que $X \to A^{1}$ induise un \uncertain{isom.}

$f^{\mathrm{fl}}_{*}(\Omega^{1}_{X/S}) \leftarrow g^{1\,\mathrm{fl}}_{*}(\Omega^{1}_{A^{1}/S})$.

\uncertain{Il suffit pour ceci que} $\Omega^{1}_{X/S}$ \struck{soit} \add{plat et} coh.\ plat en dim 0, et que les \ill{} $H^{0}(A^{1}_{s}, \Omega^{1}_{A^{1}_{s}/s}) \to H^{0}(X_{s}, \Omega^{1}_{X_{s}/s})$ \uncertain{soient} \uncertain{bijectifs}. C'est le cas en particulier si $X$ est un torseur sous schéma abélien, \ill{} si $X$ est une courbe relative \add{\ill{} $f$ \uncertain{propre et plat} \ill{}} propre et lisse. Alors $P^{c0}_{X/S}$ \add{de $(A, B)$} est une extension du schéma abélien $B$ par un fibré vectoriel, et \ill{} \uncertain{constitue} l'extension \emph{universelle} [compte tenu que
\[
B^{*} \simeq A, \quad \underline{\mathrm{Lie}}(B^{*})^{\vee} = \underline{\mathrm{Lie}}(A)^{\vee} \simeq g^{0}_{*}(\Omega^{1}_{A/S}) \simeq g^{1}_{*}(\Omega^{1}_{A^{1}/S}) \simeq f_{*}(\Omega^{1}_{A^{1}/S}) \ldots
\]
\note{après $B^{*} \simeq A$, une expression biffée \ill{} d'indice $B^{*}$ ; le dernier terme, lu $f_{*}(\Omega^{1}_{A^{1}/S})$, est peut-être $f_{*}(\Omega^{1}_{X/S})$ ; le crochet ne se ferme pas sur la page}

\page{207}
Considérons la suite exacte (1.2), où $u^{*}(P^{c}_{X/S}) \to P^{c}_{X/S}$, je dis qu'elle se factorise
\[
(2.3) \qquad u^{*}(P^{c}_{X/S}) \to P^{ci}_{X/S},
\]
\struck{En effet} dès que l'on a
\[
(2.4) \qquad f^{\mathrm{fl}}_{*}(\Omega^{1}_{X/S}) = f^{\mathrm{fl}}_{*}(\underline{Z}^{1}_{X/S})
\]
(\uncertain{identité}, \struck{\uncertain{des différentielles}} \add{\uncertain{condition}} \uncertain{dès} $f^{\mathrm{fl}}_{*}(\underline{Z}^{1}_{X/S})$ \struck{est faisceau} \add{\uncertain{noyau des}} \ill{} \uncertain{les} $f^{\mathrm{fl}}_{*}(\underline{\Omega}^{1}_{X/S}) \to f^{\mathrm{fl}}_{*}(\underline{\Omega}^{2}_{X/S})$)], et que \struck{$u^{*}(P^{c}) \to$} \struck{En effet,} \ill{}
\[
(2.5) \qquad f^{*}V = f_{*}(\underline{\Omega}^{2}_{X/S})
\]
\uncertain{est} \uncertain{représentable} par un module \uncertain{quasi-cohérent}, \uncertain{représentable} par un fibré vectoriel.
\note{devant (2.5), un symbole biffé \ill{} au-dessus de $f^{*}$}
\uncertain{En effet}, \ill{} \ill{} $\mathrm{Hom}(B, V \struck{\ill{}}) = 0$, donc
\[
\mathrm{Hom}(u^{*}(P^{c}), V) \to \mathrm{Hom}(f^{\mathrm{fl}}_{*}(\Omega^{1}_{X/S})/-, V)
\]
est injectif, \uncertain{donc} \ill{} \ill{} est \ill{} la différentielle $f^{\mathrm{fl}}_{*}(\Omega^{1}_{X/S}) \to f^{\mathrm{fl}}_{*}(\Omega^{2}_{X/S})$, qui est supposée nulle\,…

\uncertain{En particulier}, sous les conditions favorables (1.2), (2.4) et (2.5), on a
\[
(2.6) \qquad (P^{c0}) \subset P^{ci} \quad \text{i.e.} \quad P^{c0} = P^{ci0}.
\]
\marginal{un grand « 2 » dans la marge gauche, à la hauteur de (2.4)}

③. Supposons que $p.1_{S} = 0$ ($p$ un nb premier). On définit alors
\[
P^{cip}_{X/S} \subset P^{ci}_{X/S}
\]
le sous-faisceau correspondant aux modules à connexion \struck{compatibles aux puiss.\ $p$-ièmes} \struck{\ill{} faisceaux} \add{\ill{}}. Supposons que $X/S$ ait loc.\ une $p$-base, donc \ill{} :

\begin{tikzcd}[column sep=small, row sep=small, nodes={font=\scriptsize}]
(3.1) \quad 0 \arrow[r] & P^{cip}_{X/S} \arrow[r] & P^{ci}_{X/S} \arrow[r, "K"] & f^{\mathrm{fl}}_{*}(\Omega^{1\otimes p}_{X/S})
\end{tikzcd}

\marginal{au-dessus de la flèche $K$ : « $p$-courbure » ; sous elle : « \struck{\ill{}} tenseur de $p$-courbure \struck{\ill{}} = défaut de comp.\ aux puissances $p$-ièmes »}

\page{208}
\emph{Question}. Soit $E$ module à connexion \add{à courbure nulle} \struck{intégrable} sur $X/S$ en car.\ $p > 0$. Peut-on définir un tenseur \struck{\ill{}} de $p$-courbure sur $X$, qui soit
\[
K \in \Gamma(X, \underline{\mathrm{End}}(E) \otimes (\Omega^{1}_{X/S})^{\otimes p}) \ ?
\]

On a aussi une \struck{\uncertain{homomorphie}} \uncertain{suite} exacte, où $\underline{\mathrm{Pic}}^{ci}_{X/S}$ désigne l'image de $P^{ci}_{X/S} \to \underline{\mathrm{Pic}}_{X/S}$ (cf.\ (2.2) dans le cas lisse)
\[
\begin{aligned}
(3.2) \qquad 0 &\to \underbrace{[f^{\mathrm{fl}}_{*}(\underline{\mathcal{O}}^{*}_{X/S}) \to f^{\mathrm{fl}}_{*}(\underline{Z}^{1}_{X/S}) \to f^{\mathrm{fl}}_{*}(\underline{\Omega}^{1\otimes p}_{X/S})]}_{\underline{\mathrm{Cart}}_{X/S} \text{ « faisceau de Cartier »}} \to P^{cip}_{X/S} \\
&\to \underline{\mathrm{Pic}}^{ci}_{X/S} \to f^{\mathrm{fl}}_{*}(\Omega^{1\otimes p}_{X/S})/\mathrm{Im}\, f^{\mathrm{fl}}_{*}(\underline{Z}^{1}_{X/S}),
\end{aligned}
\]
\marginal{au-dessus du crochet : « (défaut d'exactitude dans le milieu) »}
où la flèche
\[
K_{0} : f^{\mathrm{fl}}_{*}(\underline{Z}^{1}_{X/S}) \to f^{\mathrm{fl}}_{*}(\underline{\Omega}^{1\otimes p}_{X/S})
\]
est donnée par
\[
(3.3) \qquad K_{0}(\omega) = \omega^{\otimes p} - F^{*}_{X/S}(\tilde{C}\omega) \quad \text{pour } d\omega = 0,
\]
\marginal{sous $\tilde{C}$ : « opération de Cartier »}

\emph{NB} $K_{0}$ \add{est additif mais} n'est pas linéaire, c'est la somme d'une opération $p$-linéaire et d'une opération linéaire.

\uncertain{Ceci} dans le cas de l'existence locale d'une $p$-base, on sait qu'il \uncertain{y a} équivalence \add{\uncertain{entre}} \uncertain{les modules} \struck{\uncertain{connexions}} \uncertain{p.\ ex.}\ $E$ sur $X$, \uncertain{munis} d'une connexion à courbure et $p$-courbure \uncertain{nulles}, i.e.\ \struck{\ill{}} sur $E$ équivalente à la \emph{donnée} d'une \emph{descente} de $E$ : $\Lambda =$ \struck{$\mathcal{O}^{p}_{X}$} $\underline{\mathcal{O}}_{S}[\underline{\mathcal{O}}_{X}^{p}] \simeq \underline{\mathcal{O}}_{X^{(p/S)}}/\mathcal{J}$, où $\mathcal{J}$ est le noyau de $\underline{\mathcal{O}}_{X^{(p/S)}} \to F_{X/S*}(\underline{\mathcal{O}}_{X})$, \struck{et qui \ill{}} \add{\ill{}} \uncertain{ou} \ill{} aux \uncertain{données} $E$ équivalent à celles des modules quasi-cohérents sur $(X^{(p/S)}, \Lambda) \simeq (X, \underline{\mathcal{O}}_{X}^{p})$. \uncertain{Sauf} $(X^{p/S\,!})$ est \uncertain{un} \uncertain{schéma}, \ill{} \uncertain{on trouve donc} que $P^{cip}_{X/S}$ est le faisceau associé au préfaisceau
\[
S' \mapsto \mathrm{Pic}(X_{S'}^{p/S'\,!}).
\]

\page{209}
Si en particulier pour tout $S'$, $X_{S'}^{p/S'} \xleftarrow{\sim} X_{S'}^{p/S'\,!}$, ce qui \uncertain{est} \uncertain{le cas} en particulier si $X/S$ est lisse, on trouve donc
\[
(3.4) \qquad P^{cip}_{X/S} \simeq \underline{\mathrm{Pic}}_{X^{(p/S)}/S} \simeq \underline{\mathrm{Pic}}_{X/S}^{(p/S)}
\]
\note{dans le deuxième terme, l'exposant $(p/S)$ de $X$ a été surchargé ; un mot au-dessus \ill{}}
L'homomorphisme canonique $P^{cip}_{X/S} \to \underline{\mathrm{Pic}}_{X/S}$ \uncertain{n'est autre} que la Verschiebung
\[
(3.5) \qquad V : \underline{\mathrm{Pic}}^{p/S}_{X/S} \to \underline{\mathrm{Pic}}_{X/S}
\]
[\uncertain{ainsi} qu'il faudrait définir $V$ avec une généralité suffisante pour que ce qui \uncertain{précède} ait \uncertain{un sens}\,…]. On trouve donc, \uncertain{grâce à} (3.2), un isom.
\[
(3.6) \qquad \mathrm{Ker}(V_{\underline{\mathrm{Pic}}_{X/S}}) \simeq
\left\{
\begin{array}{l}
\text{faisceau des sections de } f^{\mathrm{fl}}_{*}(\underline{\Omega}^{1}_{X/S}) \text{ telles} \\
\text{que } d\omega = 0,\ C\omega = \tilde{\omega}, \text{ modulo les sous-} \\
\text{faisceau des } df/f,\ f \text{ dans } f^{\mathrm{fl}}_{*}(\underline{\mathcal{O}}^{*}_{X/S})
\end{array}
\right.
\]
\note{au-dessus de « sous-faisceau », un mot biffé (\uncertain{formes}) remplacé par « \uncertain{sections} »}

Quand on \uncertain{inverse} \uncertain{sur} $(\mathrm{Sch})/S$ les \struck{\ill{}} morphismes \uncertain{radiciels} \add{[en fait, il suffit \struck{\ill{}} \uncertain{certains} \uncertain{très particuliers}]} \ill{}, \struck{\ill{}} \uncertain{de sorte} que $\underline{\mathrm{Pic}}_{X/S} \xrightarrow{F} \underline{\mathrm{Pic}}^{p/S}_{X/S}$ devient \uncertain{inversible}, \ill{} les \struck{\ill{}} premiers \uncertain{termes} \uncertain{devenant} $p^{\underline{\mathrm{Pic}}_{X/S}}$\,….

\emph{Comparaison} \struck{\emph{des}} $P^{cip\,0}_{X/S}$ et $P^{cip\,0}_{\mathrm{Alb}^{0}_{X/S}}$. \uncertain{Plaçons-nous pour} simplifier dans le cas $X/S$ lisse, \uncertain{de sorte} que (3.4) \uncertain{est valable}, et supposons que $\underline{\mathrm{Pic}}^{0}_{X/S}$ soit un schéma abélien $B$. Comme \ill{} \uncertain{aussi}
\[
(\underline{\mathrm{Pic}}_{X/S})^{(p/S)\,0} = (\underline{\mathrm{Pic}}^{0}_{X/S})^{p/S} = B^{(p/S)},
\]
\struck{\ill{}} et \uncertain{qu'on a} ainsi $B \simeq \underline{\mathrm{Pic}}^{0}_{A^{1}/S} = \underline{\mathrm{Pic}}^{0}_{A/S}$,

\page{210}
on trouve que
\[
(3.7) \qquad P^{cip\,0}_{X/S} \xleftarrow{\sim} P^{cip\,0}_{A^{1}/S} \simeq P^{cip\,0}_{A^{0}/S} \simeq B^{(p/S)}
\]
i.e.\ on trouve la \uncertain{même} \uncertain{chose} que \uncertain{dans} le cas d'un schéma \uncertain{abélien}. \add{Remarque} D'ailleurs, faisant maintenant $X = A$, on sait que $B^{cp0}$, regardé comme extension de $B$ par un groupe \struck{$\mathrm{Ker}\, V$} $\alpha_{B} = \mathrm{Ker}(V_{B} : B^{p/S} \to B)$, qui est un groupe affine sur $S$ annulé par $V$, est l'extension \emph{universelle} de cette nature [\uncertain{dès} que $S$ est le spectre d'un corps parfait]. \uncertain{Comme} $D(\alpha_{B}) = \mathrm{Ker}\, F_{B^{*}}$ \add{\struck{\ill{}} $F_{A}$} est un groupe infinitésimal de hauteur 1, dont l'algèbre de Lie est $\underline{\mathrm{Lie}}(B^{*})$, et \uncertain{l'on} sait que $\underline{\mathrm{Lie}}(B^{*})^{\vee}$ est \ill{} \emph{\uncertain{l'analogue}} \emph{vectoriel} \uncertain{des} groupes finis ; donc \ill{} \uncertain{voit} que \ill{} extensions vectorielles universelles est \uncertain{déduite}, \uncertain{par} $\alpha_{B} \to$ \uncertain{enveloppe} vectorielle $\simeq \underline{\mathrm{Lie}}(B^{*})^{\vee}$, de l'extension universelle par les groupes annulés par $V$.
\note{la phrase est très incertaine : la lecture ne garantit que les formules}

4. Supposons que $p$ soit nilpotent sur $S$, sans plus. Soit
\[
P^{cin} \subset P^{ci},
\]
défini par les connexions à courbure nulle \emph{nilpotentes}. Si $X/S$ est lisse, \struck{et} \add{les valeurs des foncteurs $P^{cin}$ sur les $S'$ \emph{réduits} sont} \struck{$S$ réduit, alors il s'agit de connexions} \uncertain{les} \uncertain{mêmes} que \struck{\ill{}} ceux de $P^{cip}$ ; \struck{\uncertain{plus généralement}} \add{\uncertain{plus généralement}}

\page{211}
\uncertain{Pour} simplifier, \uncertain{supposons} $S$ \add{loc.}\ \uncertain{noeth.}, donc \uncertain{les} $S'$ aussi, \add{\uncertain{tous}} (\uncertain{OB} \ill{}) sont \add{\ill{}} (i.e.\ \uncertain{de type fini sur $S$}) dans \uncertain{l'énoncé} \uncertain{nilpotents} \struck{\uncertain{(sur $S$)}}. On trouve que, \uncertain{une} connexion \uncertain{à courbure nulle} sur un Mod.\ inversible sur $X_{S'}$ est \uncertain{à connexion} \emph{nilpotente} ssi, en \uncertain{restreignant} \uncertain{sur} $S'_{\mathrm{red}}$, on trouve une connexion compatible aux puissances $p$-ièmes. On trouve que $P^{cin}_{X/S}$ est l'image inverse dans $P^{ci}_{X/S}$ des \uncertain{sous-faisceaux}
\[
\overline{f^{\mathrm{fl}}_{*}(\Omega^{1\otimes p}_{X/S})} \subset f^{\mathrm{fl}}_{*}(\Omega^{1\otimes p}_{X/S}),
\]
\uncertain{ce qui est} le « \uncertain{plus grand} sous-groupe $\pm$formel \uncertain{connexe} »\,…. On trouve \uncertain{ainsi}
\[
(4.1) \qquad 0 \to P^{cip} \to P^{ci}_{X/S} \to f^{\mathrm{fl}}_{*}(\Omega^{1\otimes p}_{X/S})/\overline{f^{\mathrm{fl}}_{*}(\Omega^{1\otimes p}_{X/S})}
\]
\note{l'exposant du premier terme se lit « cip » ; on attendrait $P^{cin}$}
\uncertain{et} (3.1) \uncertain{devient}
\[
(4.2) \qquad 0 \to P^{cip}_{X/S} \to P^{cin}_{X/S} \to \overline{f^{\mathrm{fl}}_{*}(\Omega^{1\otimes p}_{X/S})},
\]
\note{sous $P^{cip}_{X/S}$ : « $\simeq \underline{\mathrm{Pic}}^{p/S}_{X/S}$ »}
qui explicite $P^{cin}$ comme \uncertain{extension} d'un ss-groupe d'un groupe $\pm$formel par le groupe $P^{cip}$, qui est \uncertain{bien compris}\,…
\struck{On peut aussi se proposer d'étudier}

On conclut
\[
(4.3) \qquad 0 \to P^{cip\,0}_{X/S} \to P^{cin\,0}_{X/S} \to \overline{f^{\mathrm{fl}}_{*}(\Omega^{1\otimes p}_{X/S})}
\]
\note{sous $P^{cip\,0}_{X/S}$ : « $\simeq \underline{\mathrm{Pic}}^{0\,p/S}_{X/S}$ »}
\struck{On peut aussi \uncertain{étudier} le \uncertain{noyau} \ill{}, et}

Étudions \add{le noyau} $\mathcal{V}$ \uncertain{de} $P^{cin}_{X/S} \to \underline{\mathrm{Pic}}_{X/S}$,
\[
(4.4) \qquad 0 \to \mathcal{V} \to P^{cin}_{X/S} \to \underline{\mathrm{Pic}}_{X/S}
\]
c'est le noyau de
\[
f^{\mathrm{fl}}_{*}(\underline{Z}^{1}_{X/S})/\mathrm{Im}\, f^{\mathrm{fl}}_{*}(\mathcal{O}^{*}_{X}) \to f^{\mathrm{fl}}_{*}(\Omega^{1\otimes p}_{X/S})/\overline{f^{\mathrm{fl}}_{*}(\Omega^{1\otimes p}_{X/S})},
\]
donc il admet une \uncertain{dévissage} \struck{\ill{}} en
\[
(4.5) \qquad 0 \to \underline{\mathrm{Cart}}_{X/S} \to \mathcal{V} \to \overline{f^{\mathrm{fl}}_{*}(\Omega^{1\otimes p}_{X/S})}
\]

\page{212}
où $\underline{\mathrm{Cart}}_{X/S}$ est le faisceau de Cartier déjà envisagé dans (3.2). Regardons le \uncertain{sous-foncteur} de \struck{\ill{}} $P^{cin}_{X/S}$ image inverse de $\underline{\mathrm{Pic}}^{0}_{X/S}$ et supposons que $\underline{\mathrm{Pic}}^{0}_{X/S}$ est \uncertain{$V$-surjectif} (\uncertain{p.\ ex.}\ un schéma abélien), de sorte que $(\underline{\mathrm{Pic}}^{0}_{X/S})^{p/S} \to \underline{\mathrm{Pic}}^{0}_{X/S}$ est un épimorphisme. On trouve donc une suite exacte
\[
(4.6) \qquad 0 \to \mathcal{V} \to [P^{cin}_{X/S}]^{0} \to \underline{\mathrm{Pic}}^{0}_{X/S} \to 0
\]
\note{à l'exposant 0 du crochet : « (image inverse de $\underline{\mathrm{Pic}}^{0}_{X/S}$) » ; dans $(\underline{\mathrm{Pic}}^{0}_{X/S})^{p/S}$, un exposant biffé \ill{}}

\emph{NB} Une façon commode de dire que $P^{cin}_{X/S}$ est \uncertain{connexe}, c'est \uncertain{de dire} que $P^{cin}_{X/S}$ est \uncertain{le complété formel de} $P^{cip}_{X/S} \simeq \underline{\mathrm{Pic}}^{(p/S)}_{X/S} \subset P^{ci}$ dans $P^{ci}$. \struck{Dans le} \uncertain{En} cas favorable où $P^{ci\,0}$ est l'extension \emph{vectorielle universelle} \add{$E$} d'un schéma abélien $B$ par un groupe \uncertain{vectoriel} $V$, on trouve que \struck{$P^{cin\,0}$} \add{\uncertain{par} $P^{cin\,0}_{X/S}$} $\supset B^{p/S}$, donc $B^{p/S} \subset E$, et on \struck{trouve} \add{\uncertain{trouve}} le complété formel \struck{\uncertain{induit}} \add{\uncertain{le long}} de $B^{p/S}$. Sa connaissance est donc nettement plus précise que celle de $\overline{E}$, complété formel de $E$ le long de sa section unité, i.e.\ le groupe formel associé à $E$.

\section{$c^{1}(L)$ style connexions et opération de Cartier}
\note{titre pris sur la page 213, une chemise portant ces seuls mots de sa main, soulignés ; le mot lu « connexions » est douteux}

\page{214}
$X/S$ lisse, on considère $H^{*}_{\mathrm{DR}}(X/S) = \mathbb{H}^{*}(X, \Omega^{\bullet}_{X/S})$ et $\mathbb{H}^{*}(X, {}^{1}\Omega_{X/S})$, où ${}^{1}\Omega_{X/S} = [\Omega^{1}_{X/S} \to \Omega^{2}_{X/S} \to \cdots]$ ($\Omega^{1}$ en degré 1). On a suites spectrales aboutissant : l'une et l'autre \ill{} \ill{} la suite spectrale, donnant une \ill{} de suites exactes en bas degrés ;
\drawing{dans la marge gauche, le quadrant $(p, q)$ de la suite spectrale, les points $1, 2, \ldots$ marqués sur l'axe des $p$, avec la légende « suite spectrale pour $\mathbb{H}^{*}(X, {}^{1}\Omega_{X/S})$ »}

\begin{tikzcd}[column sep=small, row sep=small, nodes={font=\scriptsize}]
0 \arrow[r] & {}'E^{20}_{2} \arrow[r] \arrow[d, "\simeq"] & \mathbb{H}^{2}(X, {}^{1}\Omega^{*}_{X/S}) \arrow[r] \arrow[d] & {}'E^{11}_{2} \arrow[r] \arrow[d] & {}'E^{30}_{2} \arrow[d, "\simeq"] \\
{} & H^{2}(H^{0}(X, \Omega^{*})) \arrow[r] & H^{2}_{\mathrm{DR}}(X) \arrow[r] & H^{1}(H^{1}(X, \Omega^{*})) \arrow[r] & E^{30}_{2}
\end{tikzcd}

\note{au-dessus de $\mathbb{H}^{2}(X, {}^{1}\Omega^{*}_{X/S})$, « $c(\underline{L}) \in$ » ; au-dessus de ${}'E^{11}_{2}$, « $\simeq \mathrm{Ker}(H^{1}(X, \Omega^{1}_{X/S}) \to H^{1}(X, \Omega^{2}_{X/S}))$ » ; la flèche verticale sous ${}'E^{11}_{2}$ porte un mot écrit à l'envers, \uncertain{« surj »}. Au-dessus du schéma, ${}'E^{10}_{2} = Z^{1}(X, \Omega^{*}_{X/S}) \simeq \mathbb{H}^{1}(X, {}^{1}\Omega^{*}_{X/S})$, avec une flèche vers la suite exacte de la ligne du bas, qui commence plus à gauche :}
\[
\begin{aligned}
0 &\to H^{1}(H^{0}(X, \Omega^{*}_{X/S})) \to H^{1}_{\mathrm{DR}}(X/S) \to \mathrm{Ker}[H^{1}(X, \mathcal{O}_{X}) \xrightarrow{d} H^{1}(X, \Omega^{1}_{X})] \\
&\to H^{2}(H^{0}(X, \Omega^{*})) \to H^{2}_{\mathrm{DR}}(X) \to H^{1}(H^{1}(X, \Omega^{*})) \to E^{30}_{2}
\end{aligned}
\]
avec
\[
E^{10}_{2} = \frac{Z^{1}(X, \Omega^{*}_{X/S})}{\mathrm{Im}\, C^{0}(X, \Omega^{*}_{X/S})}, \quad
E^{20}_{2} = \frac{Z^{2}(X, \Omega^{*}_{X/S})}{\mathrm{Im}\, C^{1}(X, \Omega^{*}_{X/S})},
\]
\[
E^{11}_{2} = \frac{\mathrm{Ker}[H^{1}(X, \Omega^{1}) \to H^{1}(X, \Omega^{2})]}{\mathrm{Im}[H^{1}(X, \Omega^{0}) \to H^{1}(X, \Omega^{1})]},
\]
et $E^{01}_{2}$ le noyau $\mathrm{Ker}[H^{1}(X, \mathcal{O}_{X}) \to H^{1}(X, \Omega^{1}_{X})]$.

Pour que $\exists$ connexion intégrable sur $\underline{L}$, module inversible donné sur $X$, il faut et il suffit que l'obstruction dans $\mathbb{H}^{2}(X, {}^{1}\Omega^{*}_{X/S})$ soit nulle, i.e.

ⓐ \struck{\ill{}} $c^{1}_{\mathrm{Hdg}}(L) \in H^{1}(X, \Omega^{1}_{X/S})$ nulle

ⓑ deuxième obstruction dans $H^{2}(H^{0}(X, \Omega^{*})) = \dfrac{Z^{2}(X, \Omega^{*}_{X/S})}{\mathrm{Im}\, C^{1}(X, \Omega^{*}_{X/S})}$ soit nulle.

\emph{NB} En car.\ $p > 0$, l'obstruction est dans un $\mathbb{F}_{p}$-vectoriel, \uncertain{comme} tué par $p$. Donc \add{\uncertain{pour}} \ill{} \ill{} de $p\,\mathrm{Pic}(X)$ l'obstruction est nulle. On trouve, si $X$ est propre sur $S$ \add{local}, que l'obstruction est nulle \struck{dans} si $\xi$ est dans $\mathrm{Pic}^{\tau}(X/S)$ (à vérifier\,…).
\note{un double trait vertical dans la marge gauche, le long de ce NB}

Supposons l'obstruction en question nulle, \add{$S$ en car.\ $p > 0$,} alors l'indétermination pour trouver la connexion \uncertain{intégrable} est un $Z^{1}(X, \Omega^{*}_{X/S})$, i.e.\ c'est une 1-forme \uncertain{fermée} \uncertain{fermée} ($d\omega = 0$). \struck{À quelle} \struck{Donc} D'autre part, le \emph{défaut de compatibilité aux puiss.\ $p$-ièmes} est \struck{un homomorphisme} un homom.\ additif et $p$-homogène $\underline{V}_{X/S} \to \underline{\mathrm{End}}(\underline{L})$, i.e.\ un homom.\ $v : \underline{V}^{(p)}_{X/S} \to \underline{\mathrm{End}}(L)$, i.e.\ ici une section de $\underline{\Omega}^{1\,(p)}_{X/S}$. Cette section, si on remplace la connexion $\theta$ par $\theta' = \theta + \omega$, est remplacée par
\[
\begin{aligned}
\text{\struck{$\xi \mapsto$}}\ v'(\xi) &= v(\xi) + \underline{\omega(\xi)^{p} + \theta_{\xi}^{p-1}\omega(\xi) - \omega(\xi^{p})} \\
&= v(\xi) + \underbrace{\omega(\xi)^{p} - \langle C\omega, \tilde{\xi} \rangle 1_{\underline{L}}}
\end{aligned}
\]
\marginal{sous l'accolade : « opération de Cartier » ; dans la marge gauche : « \uncertain{rappelle} \ill{} \ill{} »}

\page{215}
\note{page écrite dans la largeur de la feuille et biffée tout entière de deux longs traits obliques ; on la transcrit sans marquer la biffure ligne par ligne ; la lecture est très incertaine}
La catégorie des couples \ill{} \ill{}
\[
S \xleftarrow{\varphi_{0}} U_{0} \xrightarrow{f} U \qquad (S \text{ lisse})
\]
avec disons $f$ un monomorphisme, admet des produits (: \ill{} le produit de deux couples $(U_{0}, U)$ et $(V_{0}, V)$ est $(U_{0} \times_{S} V_{0}, U \times V)$). D'autre part le foncteur $(U, U_{0}) \mapsto$ \ill{} le faisceau \uncertain{associé} sur $D_{S}$ (\ill{} sur \ill{} des sites $C_{S}$ de \ill{} $C$\,…) commute aux \uncertain{produits} (\uncertain{plus généralement} aux limites finies). Dans le cas (habituel !) où les foncteurs \ill{} $(U, U_{0}) \mapsto$ \ill{} sont \uncertain{« ind-représentables »} \ill{} $C$ à l'aide de \ill{} \uncertain{indexés} par la \uncertain{famille} des \ill{} $U_{0}$ en $U$\,…), on \uncertain{trouve} donc des foncteurs produits \uncertain{ind-représentables}, en particulier $S \times S$ \uncertain{ind-représentable} par les \ill{}
\marginal{\uncertain{Abstraction}, et \ill{} \ill{} pour cette \ill{} des produits, \ill{} « \ill{} » \ill{}}

\page{216}
\note{suite de la page 214}

\begin{tikzcd}[column sep=small, row sep=small, nodes={font=\scriptsize}]
X \arrow[r, "F_{X/S}"] \arrow[d] & X^{p/S} \arrow[r] \arrow[dl] & X \arrow[dl] \\
S \arrow[r] & S &
\end{tikzcd}

Considérons $F^{*}_{X/S}(C\omega) \in \Gamma F^{*}_{X/S}\underline{\Omega}^{1}_{X^{p/S}/S} \simeq \Gamma \underline{\Omega}^{1\,(p)}_{X/S}$, on a alors
\[
v'(\xi) = v(\xi) + \langle \omega^{\otimes p} - F^{*}_{X/S}(C\omega), \xi^{\otimes p} \rangle, \quad \text{i.e.\ \uncertain{intrinsèquement}}
\]
\[
v' = v + \underbrace{(\omega^{\otimes p} - F^{*}_{X/S}C\omega)}_{F^{*}_{X/S}(\tilde{\omega} - C\omega)}
\]
Donc la nouvelle obstruction se trouve dans
\[
\Gamma(X, \Omega^{1\,\otimes p}_{X/S})/\mathrm{Im}[Z^{1}(X, \Omega^{*}_{X/S}) \to \Gamma(X, \Omega^{1\,\otimes p}_{X/S})]
\]
l'homom.\ étant
\[
Z^{1}(X, \Omega^{*}_{X/S}) \to \Gamma(X, \Omega^{1\,\otimes p}_{X/S}), \qquad \omega \mapsto F^{*}_{X/S}(\tilde{\omega} - C\omega) = \omega^{\otimes p} - F^{*}_{X/S}(C\omega)
\]
\marginal{« 2 », en grand, dans la marge droite}

Supposons $X$ propre sur $S$, \uncertain{et} \ill{}, on \uncertain{localisant} sur $S$, un foncteur
\[
\underline{\mathrm{Pic}}^{\mathrm{conn\,int}}_{X/S} \to \underline{\mathrm{Pic}}_{X/S}
\]
\note{au-dessous, une suite exacte en partie biffée (le premier terme, biffé, \ill{}), dont reste lisible}
\[
(f_{*}\underline{Z}^{1}_{X/S}) \xrightarrow{\gamma} f_{*}(\Omega^{1\,\otimes p}_{X/S}), \qquad f_{*}\underline{Z}^{1}_{X/S} = \mathrm{Ker}[f_{*}(\Omega^{1}_{X/S}) \xrightarrow{d} f_{*}(\Omega^{2}_{X/S})]
\]
\marginal{entouré, relié à $\gamma$ : « pas linéaire ! »}
\note{plusieurs mots biffés \ill{} sous le $\underline{\mathrm{Pic}}$ (« \uncertain{connexions intégrables} »)}

Supposons que
a) $\underline{\Omega}^{1}_{X/S}$, $\underline{\Omega}^{2}_{X/S}$, $\Omega^{1\,\otimes p}_{X/S}$ sont coh.\ plats sur $S$ en dim 0,
b) \emph{\uncertain{Coker}} $d^{1}$ et $\mathrm{Im}\, d^{1}$ sont plats sur $S$, où $d^{1} : f_{*}(\underline{\Omega}^{1}_{X/S}) \to f_{*}(\underline{\Omega}^{2}_{X/S})$
[OK si $X$ de dim relative 1 sur $S$]
c) \struck{le} $\mathrm{Ker}\, \gamma$ \struck{est} plat sur $S$.

Alors le \ill{} \ill{} \ill{} \uncertain{(ci)} \uncertain{est} \uncertain{représentable} par un fibré vectoriel
\note{la page s'arrête sur cette phrase}

\page{217}
\note{page écrite dans la largeur de la feuille, barrée tout entière de deux longs traits courbes ; on la transcrit sans marquer la biffure ligne par ligne ; la lecture est très incertaine}
\uncertain{Associé} à un faisceau $F$ sur $\mathcal{D}_{S}$ \uncertain{défini}, pour $\tau = (S \xleftarrow{\varphi_{0}} T_{0} \xrightarrow{i} J)$,

\begin{tikzcd}[column sep=small, row sep=small, nodes={font=\scriptsize}]
S & U_{0} \arrow[l, "\psi_{0}"'] \arrow[d, "j"'] & T_{0} \arrow[l, dashed, "u_{0}"'] \arrow[d, "i"] \arrow[ll, bend right=30, "\varphi_{0}"'] \\
{} & U & T \arrow[l, dashed, "u"']
\end{tikzcd}

$F(\tau) =$ \uncertain{ensemble} des \uncertain{homom.}\ $u : T \to U$ tels que \uncertain{ceci} se \uncertain{factorise} par $u_{0}$, \uncertain{tels} que $\psi_{0} u_{0} = \varphi_{0}$ (\uncertain{AB} \uncertain{non nécessairement}) \ill{} \uncertain{infinitésimaux} de \ill{} [Si les \uncertain{voisinages} \ill{} $U_{0}^{(n)}$ de $U_{0}$ dans $U$ existent (\uncertain{p.\ ex.}\ \ill{}) et si $\psi_{0}$ est \uncertain{lisse}, \ill{} $U^{(n)} \to S$ \uncertain{sont} \ill{} $\Sigma$, alors le \ill{} est \uncertain{ind-représentable} \ill{}
\[
\left( S \xleftarrow{\psi_{0}} U_{0} \xrightarrow{j^{(n)}} U^{(n)} \right),
\]
] \ill{}

foncteur $F$ \uncertain{qui} \uncertain{soit} \uncertain{représentable} par \ill{} \ill{} \uncertain{grâce à} l'hypothèse \ill{}, \struck{\ill{}} \ill{} \uncertain{d'un} faisceau sur $\mathcal{D}_{S}$, \ill{} \uncertain{définie} en \ill{} \ill{} \ill{} \ill{} faisceaux $(\mathcal{D}_{S})$ \ill{} ! [\ill{}

\struck{$(C$ \ill{}$)$} $\equiv$ \ill{}

Cas $U = U_{0} = S$, $\psi_{0} = j = \mathrm{id}$, \ill{} \uncertain{produit} \ill{} \ill{} $(\mathcal{C}$ \ill{} $!!)$].

\uncertain{On} trouve le foncteur $\tau \mapsto$ \uncertain{ensemble} des $\varphi$ prolongeant $\varphi_{0}$.
\note{d'une autre encre, en haut à droite, « $\tau = (S \xleftarrow{\varphi_{0}} T_{0} \xrightarrow{i} J)$ » ; la lettre lue $J$ y est sans doute le $T$ du schéma}

\section{Connexion canonique de Gauss-Manin}
\note{titre pris sur la page 218, une chemise portant de sa main, en haut à droite, « Connexion canonique de Gauss-Manin », puis, souligné, « \uncertain{Att.}\ 2. »}

\page{219}
\[
H^{\cdot}(\mathbb{R}^{\cdot} f_{*}(\Omega^{\cdot}_{X/S})) \to \text{\struck{$\mathbb{R}$}}\ \mathcal{H}^{0}(\Omega^{\cdot}_{X/S}) \ldots
\]
\note{la formule s'interrompt ; lecture incertaine du premier terme}

\struck{Compatibilités}

\struck{3)} \add{1)} Compatibilités pour \emph{$S, T, E$ variables} \struck{\ill{}}

2) Cas des complexes de De Rham tronqués, \uncertain{filtration} à 2 \uncertain{crans} \uncertain{suites exactes} de Hodge-De Rham et \ill{} des transversalités de Griffiths et de Deligne.

3) Variantes stratifiées en car.\ 0 \struck{($S/T$ \ill{} \ill{}) par une} (cohomologie stratifiante \uncertain{au lieu} de De Rham)

Variantes cristallines en car.\ quelconque.
\note{une accolade réunit les deux lignes du 3)}

\page{220}
\struck{En général}, \uncertain{Considérons} \uncertain{alors} \uncertain{une} \uncertain{extension}
\[
\underline{\Gamma}(G)/\underline{E} \longrightarrow Z^{1}C^{\cdot}_{S/R}(\underline{g}^{*})/d^{0}(\underline{E})
\]
\note{au-dessus de la flèche, un mot biffé \ill{} ; sous $\underline{\Gamma}(G)$ : « faisceau (\struck{\ill{}}) des sections de $G$ » ; sous $Z^{1}C^{\cdot}_{S/R}(\underline{g}^{*})$ : « complexe des diff.\ rel/$R$ à coeff.\ dans $\underline{g}$ »}

D'autre part, d'après la théorie des Modules stratifiés, tout op.\ différentiel
\[
D : \underline{g} \to \mathcal{M}
\]
\struck{définit un op.\ diff.} \struck{\ill{} factorise de façon canonique en}
\marginal{entouré, dans la marge gauche : « \uncertain{spécial} à car.\ 0 ? »}
« \uncertain{à coeff.\ cts} »
\[
\underline{g} \xrightarrow{c = d^{0}} Z^{1}C_{S/R}(\underline{g}) \xrightarrow{\tilde{D}} \mathcal{M}, \qquad \tilde{D} \circ d^{0} = D
\]
($\tilde{D}$ \uncertain{est} loc.\ \uncertain{induit} \uncertain{par} \uncertain{un} \uncertain{op.\ diff.}\ \uncertain{sur} $\underline{g} \otimes \Omega^{1}_{S/R}$ ; \struck{\ill{}} $\Delta$, \struck{\ill{}} \add{\ill{}} $\Delta \circ d^{0} = D$, \uncertain{et} $\Delta$ \uncertain{restreint} \uncertain{à} $Z^{1}C_{S/R}(\underline{g})$\,…)

$D$ \uncertain{définit} \uncertain{donc} \uncertain{une} \uncertain{composée}

\begin{tikzcd}[column sep=small, row sep=small, nodes={font=\scriptsize}]
\underline{\Gamma}(G) \arrow[r, "g^{-1}dg"] \arrow[rr, bend right=30, "\tilde{\tilde{D}}"'] & Z^{1}C_{S/R}(\underline{g}) \arrow[r, "\tilde{D}"] & \mathcal{M}
\end{tikzcd}

\note{le $\tilde{\tilde{D}}$ est entouré ; la page s'arrête là, l'argument se poursuit sans doute au-delà du lot}

\end{document}
